\subsubsection{Hopcroft-Karp}

\paragraph{Idea intuitiva.}
Matching bipartito maximo en $O(\sqrt{V} E)$. Algoritmo en \textbf{fases}: en cada fase, hacer BFS para encontrar los \textbf{caminos de aumento mas cortos} desde nodos libres de $L$; luego DFS solo por esos caminos. Cada fase aumenta el matching en $\ge 1$ y toma $O(E)$. La demostracion: tras una fase, todos los augmenting paths tienen longitud $\ge$ la nueva minima; eso sucede cada $\sqrt{V}$ fases (longitud minima crece geometricamente).

\paragraph{Propiedades clave (invariantes).}
\begin{itemize}\setlength\itemsep{0pt}
	\item \texttt{dist[u]} = distancia en niveles desde un nodo libre de $L$.
	\item BFS termina cuando no hay mas camino a un nodo libre de $R$.
	\item DFS respeta los niveles (solo sigue \texttt{dist[matchV[v]] == dist[u] + 1}).
	\item Cada fase aumenta el matching en $\ge 1$ arista.
\end{itemize}

\paragraph{Codigo clave.}
El BFS de Hopcroft-Karp:
\begin{verbatim}
bool bfs() {
    queue<int> q;
    for (int v : L) {
        if (matchU[v] == -1) { dist[v] = 0; q.push(v); }
        else dist[v] = -1;
    }
    bool found = false;
    while (!q.empty()) {
        int v = q.front(); q.pop();
        for (int u : g[v]) {
            if (matchV[u] != -1 && dist[matchV[u]] == -1) {
                dist[matchV[u]] = dist[v] + 1;
                q.push(matchV[u]);
            }
            if (matchV[u] == -1) found = true;  // llegamos a libre
        }
    }
    return found;
}
\end{verbatim}

\paragraph{Complejidad.}
\begin{itemize}\setlength\itemsep{0pt}
	\item $O(\sqrt{V} E)$.
	\item Para $V = 10^4$, $E = 10^5$: $\sim 3 \cdot 10^7$ operaciones.
\end{itemize}

\paragraph{Donde tocar para modificar.}
\begin{itemize}\setlength\itemsep{0pt}
	\item \textbf{Reconstruir el matching}: ya viene en \texttt{matchU}/\texttt{matchV}.
	\item \textbf{Aniadir pesos}: Hungarian algorithm.
	\item \textbf{Densidad alta}: si $E \approx V^2$, sigue siendo util pero verificar constantes.
	\item \textbf{Aniadir capacidades}: partir los nodos con capacidad $> 1$ en varias copias.
\end{itemize}

\paragraph{Trampas y errores frecuentes.}
\begin{itemize}\setlength\itemsep{0pt}
	\item \texttt{dist[matchV[v]] == dist[u] + 1} (no \texttt{==} a secas); sin esto, el DFS no respeta los niveles y se pierde la complejidad.
	\item El BFS debe marcar los niveles correctamente; resetear \texttt{dist} para nodos alcanzables solo via matching.
	\item El DFS debe respetar los niveles; cualquier ``atajo'' invalida la cota.
\end{itemize}

\paragraph{Codigo.}
Ver \texttt{cpp.tex}, seccion \textit{Hopcroft-Karp}.
